\section{A native joint-Wilson test of the marked torus character variety}
\label{sec:torus-joint-wilson}

The preceding cut tests recover finite gluing data from which compatible
prices are computed.  Here we use a different, joint spectral test to
recover a classical configuration space.  No classical holonomy is used
to place the tensor indices. The two sites are incoming and outgoing copies of the torus,
placed at two fixed receiver sites with unit separation. The geometric input
is the standard interpretation of \(\mathrm{SU}(2)\) Wilson loops as
holonomy traces, not a prescribed metric on their spectrum.

\subsection{The physical operators and the twist convention}

Let \(r=k+2\) and \(\zeta_r=e^{\pi i/r}\).  The chiral
\(\mathrm{SU}(2)_k\) torus Hilbert space has orthonormal charge basis
\(e_n\), \(n=1,\ldots,r-1\), where \(n=2j+1\).
Turaev--Viro theory for the unitary modular category is the doubled
theory, so these chiral operators act on one factor of the physical
torus Hilbert space \(\mathcal H_k\otimes\overline{\mathcal H_k}\)
\cite{BalsamIII}.  In the positive fundamental-trace convention, the
marked meridian and longitude operators are
\begin{equation}
 D_r e_n=2\cos(\pi n/r)e_n,\qquad
 B_r e_n=e_{n-1}+e_{n+1},\qquad e_0=e_r=0.
 \label{eq:tw-torus-basic-operators}
\end{equation}
The first is the fundamental \(S\)-matrix quotient
\(S_{2n}/S_{1n}\); the second is fusion by the fundamental object.
The oriented meridional Dehn twist \(T_r\) multiplies \(e_n\) by
\(\exp\!\big(i\pi(n^2-1)/(2r)\big)\), up to one common framing phase.
The common phase cancels in the slope operator \(C_r=T_rB_rT_r^*\);
hence
\begin{equation}
 C_r e_n=\zeta_r^{\,n+1/2}e_{n+1}
          +\zeta_r^{\,-n+1/2}e_{n-1}.
 \label{eq:tw-torus-twist-operator}
\end{equation}
We fix the marking so that the slope of \(C_r\) is the product
of the two marked cycles.  Reversing the twist and the corresponding
marking just replaces the second classical eigenangle by its negative.

For clarity, these operators can be realized in the odd subspace of
\(\mathbb C[\mathbb Z/(2r)]\), with odd basis
\((\delta_n-\delta_{-n})/\sqrt2\),
\(n=1,\ldots,r-1\).  On the ambient space define
\[
 U_r\delta_n=\zeta_r^n\delta_n,\qquad
 V_r\delta_n=\delta_{n+1},\qquad J_r\delta_n=\delta_{-n}.
\]
They obey \(U_r V_r=\zeta_r V_r U_r\).  The restrictions of
\[
 U_r+U_r^*,\qquad V_r+V_r^*,\qquad
 \zeta_r^{-1/2}(U_r V_r+U_r^{-1}V_r^{-1})
\]
to the odd subspace are exactly \(D_r,B_r,C_r\).
In particular, the last formula is selfadjoint with the displayed
phase.  This is the standard torus skein representation
\cite{FrohmanGelca,CharlesMarche}.  Its product-to-sum identity follows
directly from the Weyl relation: if \(W_{(a,b)}=\zeta_r^{-ab/2}U_r^aV_r^b\)
and \(L_{\nu}=W_\nu+W_{-\nu}\), then
\begin{equation}
 L_\nu L_\mu
   =\zeta_r^{\,\det(\nu,\mu)/2}L_{\nu+\mu}
     +\zeta_r^{-\det(\nu,\mu)/2}L_{\nu-\mu}.
 \label{eq:tw-torus-product-sum}
\end{equation}
The physical three operators are the slopes \((1,0)\), \((0,1)\), and
\((1,1)\).  No commutativity at a finite level is assumed.

\subsection{A joint rank support theorem}

For a trial triple \(p=(p_1,p_2,p_3)\in\mathbb R^3\) and
a tolerance \(\delta>0\) define the positive operator and its spectral
projector
\begin{equation}
 K_{r,p}=(D_r-p_1I)^2+(B_r-p_2I)^2+(C_r-p_3I)^2,\qquad
 E_{r,p,\delta}=\boldsymbol{1}_{[0,\delta^2]}(K_{r,p}).
 \label{eq:tw-torus-joint-penalty}
\end{equation}
This is one gauge-invariant quantum spectral projection, not a product
of incompatible sharp projections.  In a unit vector, the expectation
of \(K_{r,p}\) is the sum of the squared errors of the three Wilson
operators.  The ordinary tensor of \(E_{r,p,\delta}\) across the two
marked torus ports has native two-site price \(\log_2 \operatorname{rank} E_{r,p,\delta}\)
when nonzero, and has zero spectral value when it vanishes
\cite{OpenAIPrice}.  An independent change of the positive port
separation multiplies a finite logarithmic price and does not change the
rank test.  The following statements concern this specified projector,
not the ranks of the three unsigned Wilson tensors by themselves.

\begin{theorem}[Joint-Wilson rank support and state count]
\label{thm:tw-torus-joint-wilson}
Let
\[
 X=\{(2\cos\alpha,2\cos\beta,2\cos(\alpha+\beta)):
                 (\alpha,\beta)\in(\mathbb R/2\pi\mathbb Z)^2\}.
\]
For each \(p\in X\) and each fixed \(\delta>0\), \(E_{r,p,\delta}\) has
positive rank for all sufficiently large \(r\).  For
\(p\notin X\) and \(0<\delta<\operatorname{dist}(p,X)\), it has zero rank
for all sufficiently large \(r\).  Moreover, with normalized torus Haar
measure \(d\mu\),
\begin{equation}
 \begin{aligned}
 \lim_{r\to\infty}\frac{\operatorname{rank}E_{r,p,\delta}}{r-1}
   &=\int_{(\mathbb R/2\pi\mathbb Z)^2}
     \boldsymbol{1}_{\{k_p(\alpha,\beta)\leq\delta^2\}}\,d\mu,\\
 k_p(\alpha,\beta)
   &=\big\|(2\cos\alpha,2\cos\beta,2\cos(\alpha+\beta))-p\big\|^2 .
 \end{aligned}
 \label{eq:tw-torus-count}
\end{equation}
Thus the eventually nonzero rank support at every positive tolerance
recovers \(X\), including its degenerate character points.  In the
doubled theory, using both chiral Wilson sets in the sum of penalties
gives \(X\times X\) and the analogous product-Haar count.
\end{theorem}

\begin{proof}
For spectral exclusion, work first in the full
\(\mathbb C[\mathbb Z/(2r)]\) representation.  In the quotient of bounded
matrix sequences by sequences tending to zero in operator norm, the
classes of \(U_r,V_r\) are commuting unitaries, since
\(\|U_rV_r-V_rU_r\|=|\zeta_r-1|\).
Continuous functional calculus then identifies the class of \(K_{r,p}\)
as the image of the continuous function \(k_p\) on the torus.  This
function is bounded below by \(\operatorname{dist}(p,X)^2\).
By continuity of the negative-part functional calculus, the negative
part of \(K_{r,p}-\operatorname{dist}(p,X)^2 I\) tends to zero in norm.
The same lower bound holds after restriction to the odd subspace.
This excludes sparse spectral outliers at any fixed positive distance,
a conclusion that would not follow from normalized traces alone.

For existence, every represented triple has a representative
\(\alpha\in[0,\pi]\) after joint angle inversion.  Let \(f\) be a nonzero
smooth bump of compact support in \((-1,1)\).  Choose the nearest integer
to \(r\alpha/\pi\) clipped to
\([\lfloor r^{2/3}\rfloor,r-\lfloor r^{2/3}\rfloor]\) as the charge
center \(c_r\), and let \(\sigma_r=r^{1/3}\).  In the physical odd
basis use the normalized vector
\[
 \psi_{r,n}=\mathrm{const}\cdot f\!\left(\frac{n-c_r}{\sigma_r}\right)
                 e^{-i\beta n}.
\]
The packet avoids both missing endpoints even at \(\alpha=0,\pi\).
On its support \(D_r\) differs from \(2\cos\alpha\) by
\(O(\sigma_r/r+|c_r/r-\alpha/\pi|)=O(r^{-1/3})\).
The normalized bump changes in \(\ell^2\) by \(O(\sigma_r^{-1})\)
under a unit shift, proving
\(\|(B_r-2\cos\beta)\psi_r\|=O(r^{-1/3})\).
The coefficients in \eqref{eq:tw-torus-twist-operator} applied to the
shifted packet differ from \(e^{\pm i\alpha}\) by \(O(r^{-1/3})\);
the phases of the shifted vector supply \(e^{\pm i\beta}\).
This gives
\(\|(C_r-2\cos(\alpha+\beta))\psi_r\|=O(r^{-1/3})\).
Hence \(\langle\psi_r,K_{r,p}\psi_r\rangle=O(r^{-2/3})\), and a
small eigenvalue exists below each fixed positive tolerance.

For the count, the full trace of a fixed \(U_r^aV_r^b\) vanishes for
all large \(r\) unless \(a=b=0\).  Its trace after multiplication
by \(J_r\) has absolute value at most two: a contributing basis vector
must solve \(2n\equiv-b\pmod{2r}\).  Insert \((I-J_r)/2\) when taking
the trace of a polynomial in \(D_r,B_r,C_r\).
Such a polynomial expands as a finite sum of Fourier monomials,
and commuting full \(U_r,V_r\) changes phases only by factors tending
to one.  The physical trace normalized by \(r-1\) therefore converges
to the torus Haar trace on every fixed polynomial in the classical
symbols.  Since the spectra of \(K_{r,p}\) lie in a fixed compact
interval, polynomial approximation gives weak convergence of its
empirical eigenvalue law.  The use of a discontinuous indicator in
\eqref{eq:tw-torus-count} is then legitimate: each level set of
\(k_p\) has Haar measure zero, because it is a nonconstant real-analytic
trigonometric polynomial (the term \(2\cos(2\alpha)\) from the first
square has no other term of that frequency to cancel it).
For the doubled model the two operators act on separate factors,
normalized traces factorize, and the penalty symbol is the sum of
the two classical penalties.  The endpoint construction, lower bound,
and continuity of the count apply to this sum as well.
\end{proof}

\subsection{Identification and scope}

The known flat character space for the marked torus consists of
commuting pairs in \(\mathrm{SU}(2)\) up to simultaneous conjugation.
They can be simultaneously diagonalized, with eigenangles up to joint
inversion.  Their trace triples satisfy
\[
 x^2+y^2+z^2-xyz=4,\qquad |x|,|y|,|z|\leq2.
\]
Conversely the two roots of this relation in \(z\) are
\((xy\pm\sqrt{(4-x^2)(4-y^2)})/2\), exactly the two relative-angle
choices.  For interior \(x,y\), \(z\) distinguishes the choices
since \(\sin\alpha\sin\beta\ne0\).  If either holonomy is central,
joint inversion identifies the would-be remaining ambiguity.
Thus the triple recovers the marked compact character variety everywhere.
The character variety and its quantum torus algebra are classical
ingredients; the new statement here is their recovery from the support
and state counts of the native rank probes
\eqref{eq:tw-torus-joint-penalty}--\eqref{eq:tw-torus-count}.
This is the space of allowed marked flat data, not a construction of
a nondegenerate metric from every flat pair, not an inferred curvature
radius, and not a dynamical selection of one state.

